\chapter{全卷作战}
\section{180分钟分配}
\begin{tabularx}{\textwidth}{p{23mm}p{25mm}Y}
\toprule 阶段&参考时段&动作\\\midrule
选择10题&0--38分钟&先拿能直接识别的题；单题超过4分钟仍无有效式子，标记回访。\\
填空6题&38--62分钟&草稿保留中间式；结果检查符号、参数范围、积分上下限。\\
解答6题&62--155分钟&先读完六题，优先完整拿下模型清楚的题；10分题约12分钟，12分题约15--18分钟。\\
回访与复核&155--180分钟&先补有明确路径的缺口，再查计算；最后核对答题卡题号和选择。\\\bottomrule
\end{tabularx}
\textbf{满分训练的节奏：}限时训练时先以“小题60分钟、大题95分钟、复核25分钟”为基线。每次记录识别、计算、返工分别用了多久；方法熟而计算慢，就练化简；读不出结构，就复习本册对应方法卡。
\core{\textbf{难度与时间的含义}\quad 易：一个模型或一次直接计算；中：需要换元、分支或两个知识点串联；难：抽象证明、隐藏结构或多层条件。各节时间是熟练后的训练目标，含读题和书写，不是官方难度。}
\process{读题30秒做三件事：\textbf{圈目标}（求值、证明、分类）；\textbf{圈结构}（抵消、复合、区域对称、矩阵乘积）；\textbf{写第一式}（定义、恒等式、积分微元或秩约束）。第一式写不出时，回看目标需要哪条定理的结论。}

\section{22个位置}
下表按2021--2026年六套新结构卷归纳。高数小题的知识点轮换较大；线代题位和第16、22题更集中。表中的“启动动作”用于迅速进入相应题型章节。
\small
\begin{longtable}{p{9mm}p{57mm}p{80mm}}
\toprule 题位&近六年实际考查&启动动作\\\midrule\endhead
1&积分定阶、等价逻辑、渐近线、间断点、隐函数偏导&先分清局部极限还是定义域问题；遇积分上限先计阶。\\
2&连续极值、换序、原函数拼接、参数导数、积分极值、解结构&先找分段接点、积分区域或线性算子。\\
3&变化率、导数局部性质、递推比较、积分奇偶性、特征根、隐偏导&把文字变化率写成链式法则；判断题回定义。\\
4&零点参数、复合偏导、全轴有界解、数列发散、高阶无穷小、引力&看“任意/所有”先找必要条件；物理先投影再积分。\\
5&泰勒系数、两端敛散、参数接点、可微性、换序、极值凹性&局部看最低非零阶；全局看定义和边界。\\
6&全微分反推、复合极限、参数积分最值、换序、做功、反函数&先定位求导点与中间变量；能链式就不反解。\\
7&积分定义、积分比较、拐点参数、反常命题、绝对值可导、二重积分和&比较题尽量同区间同被积结构；定义题先写极限。\\
8&惯性指数、等价相似合同、分块伴随、矩阵变换、特征符号、置换阵&标出是等价、相似还是合同，选择相应不变量。\\
9&列空间与解包含、方程相容、空间交、伴随秩、行变换、矩阵方程&把语言改写成像空间/核空间，或增广秩。\\
10&左右初等变换、向量等价、公共表示、交换与对角化、秩差、平方零&先化结构等式；不急着设所有元素。\\
11&反常积分、幂指极限、等价配参、曲率圆、反常求参、两端敛散&先判积分端点；几何先选稳定参数。\\
12&参数二阶导、隐二阶导、弧长、极值点、渐近线、抵消极限&求导先定点，极限先看抵消阶数。\\
13&隐偏导、有理积分、隐二阶导、代换方程、积分和、曲率半径&微分运算保持原变量；积分和提取步长。\\
14&双重变限、特征根、法线、高阶导、参数隐导、链式求导&画依赖关系；一层层求导，最后代点。\\
15&常微分方程、极坐标面积、差分积分、平均速度、恰当方程、均值&先辨结构：特征根/全微分/微元。\\
16&行列式系数、初等变换逆、行列式变换、向量相关、抽象通解、Gram秩&优先利用列关系与秩，避免大规模展开。\\
17&抵消极限、导数定义、几何建模、对称积分、有理积分、极坐标积分&作为10分计算题，先化结构，再算；目标12分钟。\\
18&凹凸渐近线、方程弧长、多元极值、欧拉方程、定义反推、含参积分&通常考规范过程；接点或特殊点单独处理。\\
19&弧长曲面积、二重积分、反常体积、变限体积最值、全微分极值、多元极值&先建模，再求导/积分；检查实际区域和驻点。\\
20&方程几何最值、链式求函数、二重积分、二阶链式、二重积分、无界体积&是综合转换位；第一问结果常为后一问入口。\\
21&二重积分、凸性充要、二阶中值、插值误差、割线充要、降阶方程&2022--2025连续四年为证明；2026改为降阶方程，证明与方程均需准备。\\
22&对角化、正交二次型、线性映射、解包含二次型、合同相似、低秩矩阵幂&线代综合固定落点；先用第一问确定秩/参数/结构，再完成变换。\\\bottomrule
\end{longtable}\normalsize

\section{近年串联方式}
\begin{tabularx}{\textwidth}{p{31mm}Y}
\toprule 题面结构&把哪些知识串起来\\\midrule
含参或变限积分&换元清除参数 $\to$ 牛顿--莱布尼茨 $\to$ 特殊点用定义。2020-16与2026-18可对照练。\\
极限与导数互译&先由有限极限定函数值 $\to$ 剥离已知项 $\to$ 差商。2022-17已知可导，2025-18则从连续反推可导。\\
区域交叠与对称&画交集 $\to$ 奇函数消去/交换对称 $\to$ 极坐标分角。2024-17、2025-20、2026-17。\\
微分关系求函数&链式或全微分 $\to$ 积分时保留另一变量函数 $\to$ 边界条件 $\to$ 极值。2022-20、2024-20、2025-19。\\
线代结构综合&解空间包含 $\to$ 参数与秩 $\to$ 对称矩阵或低秩分解 $\to$ 标准形/矩阵幂。2024-22、2026-22。\\\bottomrule
\end{tabularx}
\textbf{训练重点：}近六年每套都有一道独立二重积分解答题；每套第22题都是线代综合。第17题从极限、导数、积分到几何建模均有出现。对近年卷最有价值的练习是改变题面表达后仍能识别共同结构。

\section{把步骤写成得分点}
\begin{itemize}
\item 极限：写出保留到所需阶的展开和余项，最后才约去；证明可导要明确给差商极限。
\item 极值：先定义域，再驻点，再判别；$AC-B^2=0$时回原式找方向。
\item 二重积分：画域并标交点；完整写上下限和雅可比。对称化先说明区域不变。
\item 微分方程：分型、通解、初值、定义区间；除以函数之前检查被除掉的常值解。
\item 证明：先写辅助函数及端点值，再调用定理；中值点必须落在指定开区间。
\item 线代：写出核心矩阵恒等式或秩关系；特征向量与对角元素按列对应，正交变换再单位化。
\end{itemize}
